% =====================================================================
%  Section 0 -- Introduction
%  CT-S39-DRAFT-A.  Draft only.
%
%  States the recurrence, the result, and the TWO limits: the
%  identification is exact, but leaf membership is numerical and is not
%  proven.  No new claims; every figure traces to a claim id.
% =====================================================================

\section{Introduction}\label{sec:intro}

Let $b$ be a polynomial of degree $d\ge1$ with $b(k)\neq0$ for every integer
$k\ge1$, and let $Q_n$ be the denominators of the associated polynomial
continued fraction,
\[
  Q_n=b(n)Q_{n-1}+Q_{n-2},\qquad Q_{-1}=0,\quad Q_0=1 .
\]
This paper is about one such recurrence, the quadratic case
$b(n)=3n^{2}+n+1$, and about the differential equation its generating function
satisfies. The results are of two kinds, and they are established to two
different standards. Saying which is which is the purpose of this introduction.

\paragraph{Growth.} The normalised sequence $\psi_n=Q_n/\prod_{k\le n}b(k)$
converges for every such $b$, with no hypothesis on its limit, and when the
limit $\psi_\infty$ is non-zero the denominators satisfy
\[
  \frac{Q_n}{\Gamma(dn+1)}
  =K\Big(\frac{\beta_d}{d^{d}}\Big)^{n}n^{s}\bigl(1+O(1/n)\bigr),
  \qquad s=\frac{\beta_{d-1}}{\beta_d}+\frac{d-1}{2},
\]
with $K$ given in closed form by $\psi_\infty$ and the roots of $b$. The
exponential rate here is classical --- it is a corollary of the Perron--Kreuser
theorem, and \S\ref{ssec:attribution} says so before anything else in
\S\ref{sec:growth} --- and what is offered is the exponent written at the level
of the coefficients of $b$, together with the amplitude $K$ made explicit. These
statements are proved, and \S\ref{ssec:attribution} names the one classical
normalisation that is imported rather than proved.

\paragraph{The differential equation, and the leaf.} For the quadratic case the
generating function satisfies a second-order linear ODE. After an explicit
scalar reduction, the connection attached to that ODE is shown to be the
differential module of a specific member of van der Put--Saito's rank-2 family
of type $(0,0,1)$, sitting at an explicit point of the chart $U_0$ of the moduli
space $M_{t,\lambda}$. The identification is of the full operator and not merely
of its projective class; removing the projective ambiguity is a separate step,
carried out in \S\ref{sec:projective}. Painlev\'e~V is the isomonodromy
foliation of that moduli space, and the deformation vector field is regular at
the point, so exactly one leaf passes through it. Transporting along that leaf
and recomputing the trace coordinates $x_1,x_2,x_3$ of the monodromy and Stokes
data leaves them invariant to better than 90 decimal digits, at two step sizes
and on both branches of a sign pair exchanged by an explicit Galois
automorphism.

\paragraph{The two limits.} Both are stated here rather than in a closing
section, because a reader is entitled to know them before reading the argument.

\begin{enumerate}
\item \textbf{Leaf membership is numerical.} The invariance of the trace
  coordinates under transport is a high-precision numerical agreement, checked
  against a pre-registered falsifier and accompanied by controls that displace
  the point off the leaf and that substitute a deliberately wrong deformation
  field. It is evidence, and it is strong evidence, but it is not a proof, and
  nothing in this paper upgrades it to one. The identification of the operator,
  by contrast, is exact and is proved by elimination.

\item \textbf{No Painlev\'e label is asserted for this recurrence.} The
  classification question --- which row of van der Put--Saito's Table~2.1 the
  operator occupies --- is addressed only negatively. \S\ref{sec:local} shows
  that one row is \emph{excluded} for every admissible coefficient vector, by a
  Katz-invariant argument at $z=0$. A matching multiset is a necessary condition
  and not a sufficient one, so excluding a row is not the same as assigning one,
  and no row is assigned.
\end{enumerate}

\paragraph{Method, in one paragraph.} The argument is a chain of
identifications, each checked by a control that could have failed. Growth comes
from the exact recurrence
$\psi_n=\psi_{n-1}+\psi_{n-2}/(b(n-1)b(n))$, whose increments are absolutely
summable, so that convergence follows by telescoping under (A1) alone;
positivity of $b$ then supplies $\psi_\infty\neq0$ together with a two-sided
bound. Pincherle's theorem supplies the minimal-solution interpretation of that
non-vanishing, but it is not used. The local data of the ODE --- pole orders,
Katz invariants, local exponents, and the non-apparency of the two finite
singular points --- are computed exactly and matched against the hypotheses
under which van der Put--Saito construct $M_{t,\lambda}$. The point of $U_0$ is
located by matching projective normal forms, an overdetermined system of five
coefficient equations and one non-apparency condition in four unknowns; the two
located points are shown to be the complete solution set by Gr\"obner
elimination over $\Q(i\sqrt{11})$, rather than reported as what a solver
returned. The projective ambiguity is then removed by comparing the pair
$(P,Q)$ of the scalar operator directly, which a rank-one twist does not fix.
Monodromy and Stokes data are computed by numerical analytic continuation, with
the solver first validated against the closed-form Stokes multipliers of
Kummer's equation to 32 of 32 digits, and with the sectorial and orientation
conventions not assumed but selected as the unique one of sixteen candidates
meeting four independent consistency conditions.

\paragraph{Organisation.} \S\ref{sec:notation} fixes notation and the three
standing hypotheses; a reader who knows the conventions may begin at
\S\ref{sec:growth}. \S\ref{sec:growth} treats growth, and opens with the
attribution paragraph. \S\ref{sec:local} computes the local data,
\S\ref{sec:moduli} recalls the moduli space, \S\ref{sec:point} locates the
point and \S\ref{sec:projective} removes the projective ambiguity.
\S\ref{sec:monodromy} computes the monodromy and Stokes data and
\S\ref{sec:leaf} transports along the leaf. \S\ref{sec:limits} states the scope
and limits, and \S\ref{sec:open} the open questions.
